\section{METODOLOGÍA Y DISEÑO DE LA INVESTIGACIÓN}

La celda de flujo SpotON se cebó con una mezcla de priming que incluyó Flow Cell Flush, Flow Cell Tether y BSA, y la biblioteca se cargó por el puerto SpotON. La secuenciación se ejecutó en dispositivos GridION, MinION o PromethION, con adquisición de datos en tiempo real mediante el software MinKNOW, lo que habilitó el basecalling, la demultiplexación y el análisis bioinformático posterior \cite{ONT_2025,Dash_2024}. Para poliovirus, la preparación de bibliotecas siguió las indicaciones específicas del protocolo DDNS aplicables a los amplicones de VP1 \cite{Shaw_2020,Shaw_2022_protocol}.

\subsection{Procesamiento y análisis bioinformático}

La parte inicial del procesamiento de los datos de secuenciación se realizó con la plataforma EPI2ME, que ejecutó el basecalling, la demultiplexación y la generación de secuencias consenso. A partir de esas secuencias consenso, el análisis de SARS-CoV-2 se realizó con Nextclade, que asignó a cada genoma el clado de Nextstrain y el linaje del sistema PANGO correspondientes. Donde la cobertura fue suficiente se identificaron clados de la variante Ómicron y sus subvariantes, mientras que en las muestras de baja carga viral la secuencia consenso acumuló un alto porcentaje de nucleótidos indeterminados (N) que impidió asignar clado o linaje. Para RSV, la identificación del subgrupo, la selección de referencia y el análisis posterior se apoyaron en flujos automatizados de tipificación de genoma completo; mientras que para poliovirus el análisis se centró en la región VP1 para la caracterización intratípica \cite{Le_2025,Shaw_2020}. En la interpretación de las secuencias se consideró la cobertura, la profundidad por posición y el porcentaje de nucleótidos indeterminados, dado que la tasa de error de la secuenciación por nanoporos y la complejidad de la matriz ambiental exigen un preciso control de calidad \cite{Dash_2024,Child_2023}.

\subsection{Controles de calidad, controles experimentales e interpretación}

El control positivo de extracción y amplificación fueron hisopados faríngeos con carga viral conocida, que produjeron valores de Cq tempranos en la RT-qPCR. El control negativo fue un control sin molde (NTC) para detectar contaminación, y las PCR multiplex incluyeron controles positivos y negativos para verificar el desempeño de la reacción \cite{NEB_2021}. La interpretación de los resultados moleculares consideró los valores de Cq por fluoróforo, la presencia de bandas de amplificación en electroforesis y la coherencia entre réplicas. Los resultados con Cq tardíos se interpretaron con cautela, dado que la baja carga viral, la presencia de inhibidores y la degradación del ARN reducen el rendimiento de la amplificación y la secuenciación en aguas residuales \cite{Parkins_2023,MejiaCalle_2024}.

\subsection{Consideraciones éticas y de bioseguridad}

El estudio analizó muestras ambientales de aguas residuales recolectadas en el influente de la PTAR Quitumbe, sin recolección directa de datos personales ni de muestras humanas individuales con fines de identificación. El uso de hisopados faríngeos se limitó a controles positivos de laboratorio. El procesamiento de muestras se realizó conforme a prácticas de bioseguridad acordes con la manipulación de material potencialmente infeccioso \cite{WHO_2020}, mediante el uso de equipo de protección personal, cabinas de bioseguridad y la descontaminación de superficies y residuos.
